% !PRINTED 502
\weberpage{502}
%<<CONT
das wir als symbolisches Product $Q\pi_1$ darstellen. Die Elemente
von (3) sind sowohl unter einander als von den Elementen von
$Q$ verschieden, denn aus $\varkappa_1\pi_1 = \varkappa_2\pi_1$ würde $\varkappa_1 = \varkappa_2$, und aus
$\varkappa = \varkappa_1\pi_1$ würde $\pi_1 = \varkappa_1^{-1}\varkappa$ folgen, was nicht möglich ist, wenn
$\pi_1$ nicht zu $Q$ gehört. Die Elemente von $Q\pi_1$ bilden keine
Gruppe, denn sonst müsste $\varkappa_1\pi_1\varkappa_2\pi_1 = \varkappa_3\pi_1$ oder $\pi_1 = \varkappa_1^{-1}\varkappa_3\varkappa_2^{-1}$
sein, was wieder erfordern würde, dass $\pi_1$ in $Q$ enthalten ist.

Wir nennen das System $Q\pi_1$ eine \gesperrt{Nebengruppe} zu $Q$
(innerhalb $P$). Die Nebengruppe bleibt in ihrer Gesammtheit
ungeändert, wenn $\pi_1$ durch irgend ein Element $\varkappa\pi_1$ von $Q\pi_1$
ersetzt wird, weil $Q\varkappa$ mit $Q$ identisch ist.

Nennen wir überhaupt jedes System $Q\pi$ eine Nebengruppe,
so dass $Q$ selbst darunter mitgerechnet ist, so können wir zeigen,
dass zwei Nebengruppen entweder \gesperrt{ganz identisch sind}, oder
\gesperrt{kein gemeinsames Element enthalten.} Denn sind $Q\pi_1$,
$Q\pi_2$ zwei Nebengruppen, und ist $\varkappa_1\pi_1 = \varkappa_2\pi_2$, worin $\varkappa_1, \varkappa_2$ zwei
Elemente aus $Q$ sind, so folgt, wenn $\varkappa = \varkappa_2^{-1}\varkappa_1$ gesetzt wird,
% [sic] Der Druck setzt hier ein bloßes \pi; die Herleitung verlangt \pi_1. Der Druckwortlaut bleibt erhalten.
$\pi_2 = \varkappa\pi$. Da aber $\varkappa$ zu $Q$ gehört, und folglich $Q\varkappa = Q$ ist,
so folgt
\[
Q\pi_2 = Q\varkappa\pi_1 = Q\pi_1.
\]
Ist durch (2) und (3) die Gruppe $P$ noch nicht völlig erschöpft, so können wir eine weitere Nebengruppe $Q\pi_2$ bilden
und können in der Bildung der Nebengruppen so lange fortfahren, bis $P$ in lauter Nebengruppen zerlegt ist. Da jede dieser
% [sic] Der Druck gibt für die Zahl der Elemente jeder Nebengruppe \nu statt des zuvor festgelegten Gruppengrades q; quellentreu beibehalten.
Nebengruppen $\nu$ Elemente enthält, so folgt der wichtige Fundamentalsatz von \gesperrt{Cauchy}:

1. \gesperrt{Der Grad eines Theilers einer Gruppe $P$ ist
ein Theiler des Grades von $P$.}

Als specieller Fall ist darin enthalten, dass der \gesperrt{Grad
eines Elementes der Gruppe} (\S.\,153) \gesperrt{immer ein Theiler
des Grades der Gruppe ist}.

Die Zerlegung von $P$ in seine Nebengruppen können wir
(nach \gesperrt{Galois}) sehr bezeichnend durch die symbolische Gleichung
ausdrücken
\weq{4}{P = Q + Q\pi_1 + Q\pi_2 + \cdots + Q\pi_{j-1},}
wenn
\weq{5}{p = jq}
ist. Den Quotienten $j$ nennen wir den \gesperrt{Index des Theilers}
$Q$. Diese Zahl drückt keine Eigenthümlichkeit der Gruppe $Q$
an sich aus, sondern nur eine Beziehung von $Q$ zu $P$.
